\documentclass[conference]{IEEEtran}
\IEEEoverridecommandlockouts
% The preceding line is only needed to identify funding in the first footnote. If that is unneeded, please comment it out.
\usepackage{cite}
\usepackage{amsmath,amssymb,amsfonts}
\usepackage{algorithmic}
\usepackage{graphicx}
\usepackage{textcomp}
\usepackage{xcolor}
\def\BibTeX{{\rm B\kern-.05em{\sc i\kern-.025em b}\kern-.08em
    T\kern-.1667em\lower.7ex\hbox{E}\kern-.125emX}}
\begin{document}

\title{Development of a Scalable, Physics-Aware IoT Observability Platform for Precision Agriculture: An Edge-to-Cloud Approach\\
}

\author{\IEEEauthorblockN{Juan Sebastián Sutachán Correa}
\IEEEauthorblockA{\textit{Electronics Engineer} \\
\textit{Independent Research and Development}\\
Bogotá D.C., Colombia \\
jsebastiansc935@gmail.com}
}

\maketitle

%\begin{abstract}
%This document is a model and instructions for \LaTeX.
%This and the IEEEtran.cls file define the components of your paper [title, text, heads, etc.]. %*CRITICAL: Do Not Use Symbols, Special Characters, Footnotes, 
%or Math in Paper Title or Abstract.
%\end{abstract}

%\begin{IEEEkeywords}
%component, formatting, style, styling, insert
%\end{IEEEkeywords}

\section{Introduction}
The integration of High-Precision Precision Agriculture (PA) has evolved from simple periodic monitoring to a requirement for real-time observability and automated closed-loop control. However, a critical engineering gap exists between the theoretical capabilities of centralized cloud architectures and the physical realities of remote field deployments. While standard IoT models promise seamless integration, industrial environments—characterized by intermittent connectivity and harsh environmental stressors—render purely synchronous, cloud-centric models inadequate or unreliable for mission-critical tasks such as automated irrigation or nutrient management. This proposal outlines a scalable, physics-aware observability platform designed to bridge the gap between low-level hardware constraints and high-level distributed systems.

\section{Technical Problem Statement}
Current research into agricultural IoT deployments identifies three primary engineering failures that traditional hobbyist-grade or synchronous software architectures fail to address:
\begin{itemize}
    \item \textbf{Architectural Latency and Packet Failure:} Centralized cloud computing exhibits an end-to-end latency of approximately 420 ms, compared to 95 ms for edge-enabled systems. Furthermore, cloud architectures show a packet loss rate of 4.8\%, which is triple that of edge-based systems (1.6\%). In industrial irrigation, a 4.8\% failure rate in "Close Valve" commands represents a significant risk of field flooding or resource depletion.
    \item \textbf{Energy-Security Trade-offs:} Standard TLS/SSL handshakes are described as "prohibitively expensive" for low-power microcontrollers due to the execution of millions of multiplication instructions required for modular exponentiation. Transmitting a single byte on a TelosB node consumes 59.2 µJ, making the high communication overhead of standard security protocols unsustainable for long-term battery-operated or solar-powered devices.
    \item \textbf{Data Integrity and Authentication:} Most IoT-based sensor networks lack basic encryption, transmitting unencrypted packets susceptible to tampering. Beyond parity bits, industrial reliability requires robust integrity schemes, such as Hash-based Message Authentication Codes (HMAC) using SHA-256, to ensure the statistical probability of a value-flip passing verification is effectively zero.
    \item \textbf{Cloud-Edge Desynchronization:} During a cloud outage, which can exceed 48 hours in remote regions, centralized systems lose the ability to perform model updates or historical data accumulation. Without localized Atomic Control Actions, the Cyber-Physical System (CPS) node cannot maintain operational integrity.
\end{itemize}

\section{Stochastic Sensor Drift and Material Decay}

IoT sensors are constructed from physical materials subject to natural decay, causing sensor data to drift and become progressively irrelevant over time. This phenomenon is particularly prevalent in agricultural settings, where experimental data shows significant sensitivity variance in gas and moisture sensors within as little as 3-4 months post-installment.

\subsection{Mathematical Modeling of Drift}
To maintain data integrity, the platform must model drift as a stochastic process. Following a random walk model, the increment of a sensor's drift is defined by the following Gaussian distribution:

\begin{equation}
    d_{i,t} = d_{i,t-1} + \delta_{i,t}, \quad \text{where } \delta_{i,t} \sim N(0, \sigma^2)
\end{equation}

This drift is further compounded by hardware-level constraints, specifically the \textbf{Analog-to-Digital Converter (ADC)}. In industrial prototypes using the ESP32, ADC nonlinearity and a quantified resolution uncertainty of $\pm 0.01 V $ introduce residual measurement errors that degrade the overall \textbf{Signal-to-Noise Ratio (SNR)}. 

\subsection{Distributed Compensation Strategy}
The AgroSmart-IoT platform addresses these physical limitations through a hybrid, distributed computing approach:
\begin{itemize}
    \item \textbf{Edge Layer:} The ESP32-based edge nodes execute high-speed local inference (under 5 ms) and signal conditioning to match sensor outputs to ADC input ranges, maintaining Atomic Control despite physical noise.
    \item \textbf{Cloud Layer:} High-computation models, such as Variational Autoencoders (VAE-SC), are deployed via AWS Lambda/EC2 to utilize latent space distributions for estimating the drift quantity without manual intervention.
    \item \textbf{Performance:} In simulated industrial environments, this "Physics-Aware" model achieves an average recovery rate of 91\%, significantly outperforming traditional Kalman Filter (KF) techniques which are prone to error accumulation.
\end{itemize}

\section{Project Objectives}
The primary goal of the AgroSmart-IoT project is to engineer a resilient, distributed system that mitigates physical sensor degradation and network instability through a hybrid Edge-Cloud architecture.

\subsection{General Objective}

To design and implement a scalable, physics-aware IoT observability platform that ensures 99\% data integrity and autonomous operational continuity in remote agricultural environments, utilizing a Python-based FastAPI backend and AWS cloud infrastructure.

\subsection{Specific Objectives}

To achieve the general objective, the following technical milestones must be met:

\begin{enumerate}
    \item \textbf{Edge Layer Emulation (Hardware-Software Integration):} Develop a physics-aware device emulator in Python that simulates stochastic sensor drift ($d_{i,t}$) and ADC quantization noise, providing a high-fidelity data source for system validation without requiring physical deployment in the initial phase.
    \item \textbf{High-Concurrency Ingestion Layer (Backend Engineering):} Architect an asynchronous API using \textbf{FastAPI} and \textbf{Uvicorn} capable of handling up to 100 concurrent simulated sensor nodes with a 95th percentile latency ($\text{P95}$) of less than 200 ms.

    \item \textbf{Resilient Messaging \& Security (Protocols):} Implement a telemetry pipeline using the \textbf{MQTT} protocol and \textbf{HMAC-SHA256} authentication to ensure data provenance and integrity, minimizing the energy overhead compared to standard TLS handshakes.

    \item \textbf{Cloud Infrastructure \& Orchestration (Cloud Architecture):} Deploy a serverless data processing pipeline on \textbf{AWS} (using \textbf{IoT Core}, \textbf{Lambda}, and \textbf{DynamoDB}) to perform long-term drift analysis and state persistence, ensuring system recoverability after 48-hour network outages.

    \item \textbf{System Validation (Observability):} Evaluate the platform's performance by comparing the recovery rate of the Variational Autoencoder (VAE) model against traditional baseline filters, aiming for a 90\%+ accuracy in drift compensation.
\end{enumerate}

%\section{System Architecture and Design}


%\section{Implementation Methodology}


%\section{Verification and Validation}


%\section{Results and Discussion}


%\section{Conclusion}


%
%\begin{table}[htbp]
%\caption{Table Type Styles}
%\begin{center}
%\begin{tabular}{|c|c|c|c|}
%\hline
%\textbf{Table}&\multicolumn{3}{|c|}{\textbf{Table Column Head}} \\
%\cline{2-4} 
%\textbf{Head} & \textbf{\textit{Table column subhead}}& \textbf{\textit{Subhead}}& \textbf{\textit{Subhead}} \\
%\hline
%copy& More table copy$^{\mathrm{a}}$& &  \\
%\hline
%\multicolumn{4}{l}{$^{\mathrm{a}}$Sample of a Table footnote.}
%\end{tabular}
%\label{tab1}
%\end{center}
%\end{table}

%\begin{figure}[htbp]
%\centerline{\includegraphics{fig1.png}}
%\caption{Example of a figure caption.}
%\label{fig}
%\end{figure}


\bibliographystyle{IEEEtran}
\bibliography{references}


\end{document}
